\chapter{Probabilités \& Information}

L'IA ne modélise pas la certitude absolue, mais navigue dans l'incertitude. La théorie des probabilités offre le cadre mathématique pour gérer cette incertitude, tandis que la théorie de l'information quantifie les données.

\section{Distributions de Probabilités}

\begin{definition}[Distribution Normale]
La loi normale (ou Gaussienne) de moyenne $\mu$ et d'écart-type $\sigma$ a pour densité de probabilité :
$$ f(x) = \frac{1}{\sigma\sqrt{2\pi}} e^{-\frac{1}{2}\left(\frac{x-\mu}{\sigma}\right)^2} $$
\end{definition}

\section{Théorie de l'Information}

\begin{definition}[Entropie de Shannon]
L'entropie d'une variable aléatoire discrète $X$ ayant une distribution de probabilité $P$ est définie par :
$$ H(X) = - \sum_{x \in X} P(x) \log P(x) $$
Elle mesure la quantité d'incertitude ou de "surprise" inhérente aux résultats possibles de $X$.
\end{definition}

\section{Entropie Croisée}

L'entropie croisée est la fonction de perte la plus courante pour les problèmes de classification en apprentissage automatique.

\begin{definition}[Entropie Croisée]
L'entropie croisée entre une vraie distribution $P$ et une distribution estimée $Q$ est :
$$ H(P, Q) = - \sum_{x \in X} P(x) \log Q(x) $$
Minimiser l'entropie croisée équivaut à minimiser la divergence de Kullback-Leibler (KL) entre $P$ et $Q$.
\end{definition}
